\scanpage{78}
Pour qu'un sous-fibré vectoriel \inserted{$L$} de $\mathfrak g$ de corang $1$ soit de la forme $\underline{\operatorname{Lie}}(B)$, il f. et s. que le fibré \deleted{$\unsure{L}$} \inserted{$L^\perp$} de rang $1$ orthogonal \deleted{$\unsure{L^\perp}$} dans $\widetilde{\mathfrak g}$ satisfasse
\margin{\uline{Question} suffit-il que $L$ soit une ss-Algèbre de Lie i.e. stable par crochet ??}
% slanted
\[
(q\mid L^\perp)=0
\]
\margin{\uline{NB} Il \unsure{revient|reprends} au $=$ de dire que $L^\perp$ \deleted{\unsure{soit}} \inserted{est} formé d'él. de \uline{carré nul} de $M$}
% slanted
(et ceci, bien sûr, caractérise $q$ \inserted{sur $\widetilde{\mathfrak g}$ (défini comme $\mathfrak g^\vee$)} à scalaire inversible près \inserted{en termes de $G$} -- donc au signe près si on veut qu'il soit défini sur tte base de façon compatible avec chgt de base)

\deleted{\unsure{Hyp|NB}} \uline{NB} Plaçons-nous \uline{en car $2$} Alors
\[
\mathcal O_S1\subset\widetilde{\mathfrak g}\tag*{(58)}
\]
et le \deleted{\unsure{l'application|application}} morphisme
\[
X\longrightarrow\deleted{\unsure{\mathrm{Gr}}}\unsure{P|\mathbb P}(\widetilde{\mathfrak g})\tag*{(59)}
\]
allant des Borels de $\widetilde{\mathfrak g}$ vers leurs algèbres de Lie, \unsure{considéré|considérés|considérées} comme sous fibré vect. de corang $1$, est celui qui à tte section de $X$, identifiée à une \inserted{fibré en} \unsure{droites|droite} \inserted{$L$} dans $\widetilde{\mathfrak g}$ satisfaisant $q\mid L=0$, associe son orthogonal dans $\widetilde{\mathfrak g}$ pour $\varphi$. \deleted{L'image est donc les} \inserted{Comme} l'orthogonal contient $\mathcal O_S.1$, (qui est le noyau de $\varphi$ \unsure{\}|,} i.e. de $f$), \deleted{est} le morphisme précédent se factorise par
% archive sheet number: 77
